%% Numbers in this section come from generated/numbers.tex and
%% generated/table_*.tex (built by tools/make_assets.py from the fetched
%% experiment artifacts). Do not type results here by hand.
\section{Evaluation}
\label{eval}

We evaluate the volume-set MTPP on three axes that mirror its intended
role in TFOW: (i)~one-step predictive quality of the next event set, its
timing, and its type, against adapted neural and non-parametric baselines
(Sec.~\ref{sec:eval-set}--\ref{sec:eval-timetype}); (ii)~closed-loop
stability of the learned dynamics and generative realism of free-running
simulations under the stylized facts of high-frequency order flow
(Sec.~\ref{sec:eval-stylized}); and (iii)~price-level realism of
simulated books against recorded ground truth
(Sec.~\ref{sec:eval-price}).  Together these establish the model as a
faithful generative model of LOB microstructure---the property the 3S
detection and query-localization stages rely on.

\subsection{Evaluation Setup}\label{sec:eval-setup}

\paragraph{Data.}
We use Level-2 order book and trade data for BTC/USDT, ETH/USDT and
SOL/USDT on Binance (Kaiko feed), reconstructed into atomic events by
the event constructor of Section~3 (MO/CO/LO/IS $\times$ side $\times$
top $K{=}10$ levels; 62 fixed event types).  The corpus covers seven
consecutive trading days (1--7 January 2026) for each of the three
pairs (21 daily files).  Sliding windows of 50 events (stride 32) are
split chronologically \emph{within each asset} 70/15/15 into
train/validation/test, with a one-window gap at each boundary to
prevent leakage.  Timestamps are exchange-stamped on a $0.1$\,s grid;
the mean target set contains \dataTrueSetSize{} of 62 types,
confirming that simultaneous events are pervasive.

\paragraph{Models.}
We train the volume-set MTPP with three interchangeable continuous-time
decoders: a Neural-Hawkes cell~\cite{Mei2017TheProcess}
(\textsc{vsm-hawkes}), an RMTPP-style recurrent
decoder~\cite{du2016recurrent} (\textsc{vsm-rmtpp}), and a structured
state-space decoder (\textsc{vsm-s2p2}).  Baselines follow the
borrowed-backbone protocol of Shi and
Cartlidge~\cite{Shi2022StateSimulation}: an LSTM, a causal Transformer
(SAHP-style~\cite{zhang2020self}), a decayed-cell recurrent network
(CT-LSTM-style~\cite{Mei2017TheProcess}), and a per-type parallel
recurrent network (PCT-LSTM-style~\cite{Shi2022StateSimulation}), each
fitted with identical set/volume/time heads, multi-hot event-set
histories, and the same training budget (30 epochs, batch 512, AdamW,
$10^{-3}$).  Our models carry an explicit per-channel volume head
(log-volume Normal) whose negative log-likelihood over active channels
enters the joint objective.\footnote{These are \emph{adapted} backbones
under a shared set-decoding protocol, not paper-faithful
continuous-time likelihood implementations; comparisons are against the
architectures' inductive biases, not the original systems.}  A
non-parametric empirical transition baseline (\textsc{EmpTrans};
next-set frequencies conditioned on the last event set) anchors the
tables.  Results are from a single seed; a three-seed extension is in
progress.

\paragraph{Protocol.}
All headline metrics are \emph{target-free}: no model receives the
ground-truth next-event time.  For the volume-set MTPP, the predictive
distribution over the next arrival is obtained from the learned ground
intensity by trapezoid quadrature of the compensator
$\Lambda(t)=\int_0^{t}\lambda(s)\,ds$ on a log-spaced grid; point
predictions read out the decision-theoretically matched statistic
(predictive median for MAE, predictive mean for RMSE), and event-set
probabilities are evaluated at the model-predicted time.  Set decoding
uses either a global threshold calibrated for validation
F\textsubscript{1} (shared with all baselines) or a per-sample
expected-F\textsubscript{1} rule, $k^{*}=\arg\max_k 2\sum_{i\le
k}p_{(i)}/(k+\sum_j p_j)$~\cite{ye2012optimizing}.  As a control we
also report the time-conditioned (``known-time'') variant; the
model's advantage does not derive from timing information.

\subsection{Event-Set Prediction}\label{sec:eval-set}

Table~\ref{tab:set} compares set-valued next-event prediction across
all models under both decodes.  The state-space variant
(\textsc{vsm-s2p2} with expected-F\textsubscript{1} decoding) leads
every model on both per-sample Jaccard (\resSspJaccardEf{}, an $8\%$
relative improvement over the strongest adapted baseline) and
F\textsubscript{1} (\resSspFoneEf), at roughly half the predicted set
size of the neural baselines; the non-parametric \textsc{EmpTrans}
trails all learned models on set metrics.  It also attains the best
event NLL (\resSspNll{} vs.\ $\geq$\resHawkesNll{} for the other
variants and $\approx$3.10 for the baselines) and is the first learned
model in this comparison to reach the empirical transition table's
single-type accuracy (\resSspTypeAcc{} vs.\ \resEmptransTypeAcc).

\begin{table}[t]
\caption{Event-set prediction on the test set (single seed; true mean
set size \dataTrueSetSize).  Threshold decode is shared by all models
($\tau$ calibrated on validation); EF\textsubscript{1} is the
per-sample expected-F\textsubscript{1} decode.  Best in \textbf{bold}.}
\label{tab:set}
\small
%% AUTO-GENERATED by tools/make_assets.py -- do not edit.
\begin{tabular}{lccc}
\toprule
Model & Jaccard$\uparrow$ & F$_1\uparrow$ & $|\hat S|$ \\
\midrule
\multicolumn{4}{l}{\emph{Expected-F$_1$ decode}}\\
\textsc{vsm-hawkes} & 0.183 & 0.289 & 7.4 \\
\textsc{vsm-rmtpp} & 0.181 & 0.286 & 7.8 \\
\textsc{vsm-s2p2} & \textbf{0.202} & \textbf{0.307} & 7.3 \\
\midrule
\multicolumn{4}{l}{\emph{Shared global-threshold decode}}\\
\textsc{vsm-hawkes} & 0.178 & 0.281 & 8.3 \\
\textsc{vsm-rmtpp} & 0.177 & 0.278 & 8.8 \\
\textsc{vsm-s2p2} & \textbf{0.194} & \textbf{0.298} & 8.9 \\
Causal Transformer (SAHP-ad.) & 0.185 & 0.286 & 9.3 \\
Decayed cell (CT-LSTM-ad.) & 0.187 & 0.288 & 9.9 \\
PCT-LSTM (lite) & 0.174 & 0.271 & 10.3 \\
LSTM & 0.182 & 0.285 & 11.1 \\
\textsc{EmpTrans} & 0.165 & 0.265 & 6.0 \\
\bottomrule
\end{tabular}

\end{table}

\subsection{Next-Event Time, Type, and Volume}\label{sec:eval-timetype}

Table~\ref{tab:timetype} reports target-free next-event timing, the
single-next-event adaptation of Shi and
Cartlidge~\cite{Shi2022StateSimulation} (primary label $=$ a
deterministic representative of the target set; event-type NLL from
normalized channel intensities), and per-channel volume MAE (log1p
space, true channels) from the joint volume head.
\textsc{vsm-rmtpp} is the best-calibrated timer of all models: its
target-free predictive median ($0.098$\,s) essentially coincides with
the true median inter-arrival ($0.100$\,s) and its MAE
(\resRmtppTimeMae) matches the best baseline.  The
\textsc{hawkes} and \textsc{s2p2} decoders are systematically
\emph{mis}-calibrated between events (predicted intensity
$\approx$3$\times$ the true rate; Sec.~\ref{sec:eval-discussion}),
while volume MAE is statistically tied across our variants and the
best dedicated baseline head.
One honest caveat: the inter-arrival distribution is nearly degenerate
(the $0.1$\,s exchange grid concentrates much of the mass at exactly
$0.100$), so absolute timing differences are small.

\begin{table}[t]
\caption{Target-free next-event time, Shi-style single-event metrics,
and volume MAE (single seed).  Our models read out the predictive
median for MAE and predictive mean for RMSE.  \textsc{EmpTrans} NLL is
omitted (non-comparable score family).}
\label{tab:timetype}
\footnotesize
\setlength{\tabcolsep}{3.2pt}
%% AUTO-GENERATED by tools/make_assets.py -- do not edit.
\begin{tabular}{lccccc}
\toprule
Model & MAE$\downarrow$ & RMSE$\downarrow$ & NLL$\downarrow$ & Acc.$\uparrow$ & Vol.MAE$\downarrow$ \\
\midrule
\textsc{vsm-hawkes} & 0.108 & 0.199 & 2.96 & 0.202 & 0.720 \\
\textsc{vsm-rmtpp} & \textbf{0.059} & 0.194 & 2.97 & 0.221 & 0.729 \\
\textsc{vsm-s2p2} & 0.115 & 0.200 & 2.89 & 0.243 & 0.744 \\
Causal Transformer (SAHP-ad.) & 0.063 & 0.193 & 3.10 & 0.135 & --- \\
Decayed cell (CT-LSTM-ad.) & 0.065 & \textbf{0.192} & 3.10 & 0.148 & --- \\
PCT-LSTM (lite) & 0.064 & \textbf{0.192} & 3.12 & 0.153 & --- \\
LSTM & \textbf{0.059} & 0.193 & 3.11 & 0.147 & --- \\
\textsc{EmpTrans} & 0.068 & 0.197 & --- & \textbf{0.248} & --- \\
\bottomrule
\end{tabular}

\end{table}

%% TODO(tsne): planned t-SNE of the 62 learned type embeddings (4 types
%% highlighted) -- not yet produced.

\subsection{Stylized Facts of Simulated Order Flow}\label{sec:eval-stylized}

A generative model is only a \emph{simulator} if it remains
on-distribution when its own outputs are fed back as history.  We
evaluate this in closed loop: event sets, inter-arrival times, and
per-channel volumes are sampled from the model (times by inversion of
the compensator, sets from the conditional set distribution, volumes
from the log-volume head) and appended to the conditioning history,
for a fixed simulated \emph{duration} (600\,s) so that decoders with
different event rates are compared over equal market time.  Rollouts
neither explode nor die out: simulated inter-arrival and set-size
marginals remain at a constant distance from the real marginals across
the full horizon.

Following the realism-evaluation practice of market
simulators~\cite{vyetrenko2020getreal}, we aggregate real and simulated
streams into 1-second buckets and compute the eleven stylized facts of
Cont~\cite{cont2001empirical} on a signed order-flow proxy for returns
(price-moving events: market orders and best-level in-spread
insertions) and an activity series (events per bucket).
Figure~\ref{fig:stylized} shows the panel for \textsc{vsm-rmtpp}.
The decoders separate sharply.  \textsc{vsm-rmtpp} is the best
all-round simulator: near-zero return autocorrelation
(\sfRmtppAcfAbsSim{} vs.\ \sfRmtppAcfAbsReal{} real), the closest
tails (kurtosis \sfRmtppSkewReal-level), correct skew sign, a matching
activity--volatility coupling (\sfRmtppActVolSim{} vs.\
\sfRmtppActVolReal{} real), and the only non-degenerate volatility
clustering signal (ACF$|r|$ \sfRmtppAcfAbsSim{} vs.\
\sfRmtppAcfAbsReal{} real).  \textsc{vsm-hawkes} under joint volume
training simulates almost no volatility memory
(\sfHawkesAcfAbsSim); an ablation without the volume head restores
its clustering to $0.051$ and its kurtosis to near-real ($1.89$ vs.\
$1.93$), isolating a multi-task interference effect: the volume
likelihood's gradients, flowing through the shared recurrent cell,
dampen the \emph{variability} of the self-exciting intensity while
leaving prediction metrics unchanged.  \textsc{vsm-s2p2} is
degenerate in closed loop (return autocorrelation $0.89$: a
near-deterministic drift), examined next.

\begin{figure}[t]
\centering
\autofig{figs/stylized_facts_rmtpp.png}{\linewidth}
\caption{Cont's eleven stylized facts, real vs.\ simulated order flow
for \textsc{vsm-rmtpp} (single seed, 600\,s closed-loop rollouts).}
\label{fig:stylized}
\end{figure}

\subsection{Price-Level Realism}\label{sec:eval-price}

\paragraph{From events to prices.}
Simulated streams are converted to prices by a deterministic
book-replay engine that inverts the event constructor's semantics:
market orders consume the touch and walk the ladder; cancellations
that empty the best level move the quote away; in-spread insertions
improve it by their tick distance.  Events within one timestamp are
applied in matching-engine order (MO, IS, CO, LO).  The engine assumes
a dense one-tick ladder at the top of book; following the
tick-size-regime analysis of Jain et al.~\cite{Jain2025NoBooks}, this
is accurate for large-tick regimes and an approximation for small-tick
assets, so we convert ladder units to prices via recorded per-level
gaps and restrict price-level claims to distributional statements.
The engine contains no stochastic or fitted components: unlike prior
work that attaches order prices and volumes to simulated events from
separately fitted power laws with hand-tuned aggression
percentages~\cite{Shi2022StateSimulation}, every stochastic
quantity---type, side, book level, aggression, timing, volume---is
produced by the model, and only deterministic mechanics are imposed.

\paragraph{Recorded ground truth.}
The real side of every comparison uses the \emph{recorded} book state
emitted by our event constructor (mid-price and full ten-level ladders
per event), not a reconstruction.  Simulated books are initialized
from the recorded snapshot at the seed point.  Model events that are
inconsistent with the simulated book (cancellations of empty levels,
insertions into a one-tick spread) are skipped and \emph{counted}; we
report this invalid-event rate as a measure of the model's structural
coherence (\textsc{vsm-hawkes} \pvHawkesInvalid\%,
\textsc{vsm-rmtpp} \pvRmtppInvalid\%, \textsc{vsm-s2p2}
\pvSspInvalid\%).

\paragraph{Results.}
Against recorded reality (mean spread \pvHawkesSpreadRealMean\,bps),
the simulated books separate sharply: \textsc{vsm-rmtpp} is the most
faithful (\pvRmtppSpreadSimMean\,bps, with zero impossible in-spread
insertions), \textsc{vsm-hawkes} is over-tight and over-calm
(\pvHawkesSpreadSimMean\,bps), and \textsc{vsm-s2p2}'s book
\emph{diverges} to \pvSspSpreadSimMean\,bps mean
(\pvSspSpreadSimPnf\,bps at p95) --- $40\times$ reality.  Thus the
decoder that wins every one-step prediction column produces the least
viable market: one-step predictive quality does \emph{not} predict
closed-loop viability, a compounding-error phenomenon familiar from
model-based reinforcement learning that we propose as a first-class
evaluation criterion --- \emph{closed-loop book stability} --- for
LOB generative models.

\begin{figure}[t]
\centering
\autofig{figs/price_v2_hawkes.png}{\linewidth}
\caption{Price-level realism for \textsc{vsm-hawkes} against recorded
book state: mid trajectories, return Q-Q, autocorrelations, volatility
clustering, spread, signature plot, book shape, and touch-queue
density.}
\label{fig:pricev2}
\end{figure}

\subsection{Discussion}\label{sec:eval-discussion}

Four findings shape how TFOW should be used.  (1)~\emph{Set-valued
modeling is the source of predictive advantage}: the state-space
variant beats every adapted baseline on set metrics and matches the
frequency table on single-type accuracy --- multi-label structure,
not next-type guessing, is where learnable signal lives.
(2)~\emph{No decoder dominates; choose by task}: the state-space
decoder for prediction and likelihood scoring, the constant-hazard
recurrent decoder for timing and simulation --- its rigid hazard is
simultaneously why it cannot win the prediction table and why it is
the only decoder whose closed-loop market remains faithful.
(3)~\emph{Joint volume modeling is prediction-neutral but dampens
self-exciting dynamics}: volume accuracy matches dedicated baseline
heads with no cost to set/type/timing metrics, yet its gradients,
flowing through the shared recurrent state, suppress the Hawkes
decoder's closed-loop volatility clustering (restored by a
volume-head-free ablation); we report this multi-task interference as
a finding, and a stop-gradient volume readout as the remedy.
(4)~\emph{Known limitations}: the feed's $0.1$\,s grid makes
continuous hazards mis-specified (the cause of both temporal PIT
over-dispersion and the expressive decoders' between-event intensity
inflation; a grid-aware discrete-time hazard is the principled fix),
the conditionally independent set head caps within-tick co-occurrence
structure, and the study covers a single venue, one week, and three
symbols; scaling the corpus, a discrete-time temporal head, and a
within-tick dependent set head are the subject of ongoing work.
